\documentclass[10pt,a4paper]{amsart}
\usepackage[margin=19mm]{geometry}
\usepackage{amsmath,amssymb,mathtools}
\usepackage[scr=rsfs,cal=euler]{mathalfa}
\usepackage{xfrac}
\usepackage[unicode,hidelinks,bookmarks=false]{hyperref}
\newcommand{\ie}{i.e.\ }
\newcommand{\cf}{cf.\ }
\newcommand{\resp}{respectively\ }
\newcommand{\Subsection}{\subsection}
\providecommand{\nobreakdash}{\nobreak\hbox{-}}
\setlength{\emergencystretch}{2em}
\multlinegap0pt
\usepackage{bm}
\usepackage{bbm}
\usepackage{dsfont}
\usepackage{xspace}
\usepackage{cancel}
\usepackage{mathtools}
\usepackage{amsthm}
\makeatletter
\@ifundefined{theorem}{\newtheorem{theorem}{Theorem}}{}
\@ifundefined{lem}{\newtheorem{lem}[theorem]{Lemma}}{}
\@ifundefined{proposition}{\newtheorem{proposition}[theorem]{Proposition}}{}
\makeatother
\title[Nonresonance for problems involving $(p,q)$-Laplacian equati]{Nonresonance for problems involving $(p,q)$-Laplacian equations with nonlinear perturbations}
\author{Emer Lopera, Nsoki Mavinga, Diana Sanchez}
\date{Source: arXiv:2507.08229v1 (CC BY 4.0)}
\begin{document}
\maketitle

\noindent\emph{Source and licence.} Nonresonance for problems involving $(p,q)$-Laplacian equations with nonlinear perturbations. Version arXiv:2507.08229v1 (CC BY 4.0), distributed under the Creative Commons Attribution 4.0 International licence, as recorded in the arXiv licence metadata for this exact version and shown on the abstract page https://arxiv.org/abs/2507.08229v1. Full rights notice: licenses/arxiv-2507.08229-CC-BY-4.0.txt.\par\smallskip

\noindent\textit{Editorial scope.} Counted unit: the Lem whose source label is \texttt{2-exist-eigen} in the article (source0.tex lines 587--601, source label \texttt{2-exist-eigen}), delivered together with the article's own complete original proof of it (source0.tex lines 602--718, one \texttt{proof} environment of 117 lines). No mathematics is written, repaired or re-derived: statement and proof are the article's own text line for line. Required context retained uncounted: Lambda\_q (lines 161--163); prop-cu (lines 353--362); bound-K\_q (lines 497--499); def\_Nehari (lines 579--581); J\_rest\_eta (lines 584--586). Every definition, display and label of those lines is retained unchanged. Omitted from this package and not redistributed: 1--160, 164--352, 363--496, 500--578, 582--583, 719--1126. Nothing inside a retained excerpt is omitted or altered: every retained body is delivered exactly as the article's lines are written, so no omission receipt of any kind accompanies this package. Citations: the retained excerpts carry no citation command at all, so no bibliography entry of the article is transcribed and no reference is redistributed. Reference adaptation: none was needed and none was made; every reference inside the delivered unit resolves inside it, so no unresolved reference or citation is produced. Printed slips kept literal: no printed word, symbol, hypothesis or number of the counted statement or of its proof was altered. Layout and class: the article's own class is \texttt{sn-jnl} with packages including graphicx, multirow, amsmath, amssymb, amsfonts, amsthm, mathrsfs, appendix, xcolor, textcomp, manyfoot, booktabs, algorithm, algorithmicx; the wrapper is the lane's pinned \texttt{amsart} editorial wrapper at 10pt on A4, which restates the theorem-like environments the retained text needs and adds the article's own macro definitions that the retained text needs (none), with extra wrapper packages (bm, bbm, dsfont, xspace, cancel, mathtools). The delivered numbering is the unit's own. Assets: no retained span contains a figure environment, a graphics-inclusion command, a PDF-page inclusion, a table or a TikZ picture; no image file is copied into this package, the package declares no asset dependency and no image file is redistributed with it. Licence: version arXiv:2507.08229v1 of this article is distributed under the Creative Commons Attribution 4.0 International licence, as recorded in the arXiv licence metadata for this exact version and shown on the abstract page https://arxiv.org/abs/2507.08229v1. A full rights notice is delivered at licenses/arxiv-2507.08229-CC-BY-4.0.txt. Characters: every retained excerpt is pure ASCII, so the delivered engine, XeLaTeX with its default Latin Modern text font, reports no missing character.
\begin{equation}\label{Lambda_q}
\Lambda_q:= \Lambda_q(m,\rho)=\inf_{w \in \mathcal{C} \setminus \{0\}} \frac{\|w\|_{c_2,q}^q}{\int_\Omega m(x)|w|^{q}+ \int_{\partial\Omega} \rho(x)|w|^{q}},
\end{equation}

\begin{proposition}\label{prop-cu}
Let $E$ a measurable set in $\mathbb{R}^N$, $c\in L^\infty(E)$ and $p>1$. Consider the functional $I_2:L^{p}(E)\to \mathbb{R}$ defined by
\begin{equation*}
I_2(u)=\frac{1}{p}\int_{E} c(x)|u|^p.
\end{equation*}
Then, $I_2\in C^1(L^{p}(E))$ and 
 \begin{equation*}
I_2'(u)v=\int_{E} c(x)|u|^{p-2} u\cdot v  \qquad \forall  u, v \in L^{p}(E) \,.
\end{equation*}
\end{proposition}

    \begin{equation}\label{bound-K_q}
        K_q(w) \leqslant \|m\|_\infty \|w\|_q^q+ \|\rho \|_\infty \|w\|_{\partial\Omega, q}^q \leqslant C_1 \|w\|_q^q, \mbox{ for all } w\in W,
    \end{equation}

\begin{align}\label{def_Nehari}\mathcal{N}_{\lambda}&=\{v\in W\backslash\{0\}\;|\;\langle J'_\lambda (v),v\rangle=0\}\nonumber\\
    &=\{v\in W\backslash\{0\}\;|\; \|v \|_{c_1,p}^p+\|v \|_{c_2,q}^q=\lambda K_q(v)\}.
\end{align}

\begin{equation}\label{J_rest_eta}
    J_{\lambda}(v)=\frac{q-p}{pq}\|v \|_{c_1,p}^p.
\end{equation}

\begin{lem}\label{2-exist-eigen}
     If $q>p$ and $\lambda>\Lambda_q$, then
     \begin{itemize}
         \item[(i)] $ \mathcal{N}_{\lambda}\neq \emptyset.$
         \item[(ii)] Every minimizing sequence  for $J_{\lambda}$ restricted to $\mathcal{N}_{\lambda}$ is bounded in $W$, that is, if $\left(u_{n}\right) \subseteq \mathcal{N}_{\lambda}$ satisfies 
\begin{equation}\label{Jlambda}
\lim_{n\to \infty} J_\lambda (u_n) = m_{\lambda} :=\inf _{w \in \mathcal{N}_{\lambda}} J_{\lambda}(w), 
\end{equation}
then $(u_n)$ is bounded in $W$.
         \item[(iii)]  $m_{\lambda}>0$.
         \item[(iv)] 
     There exists $u_{*} \in \mathcal{N}_{\lambda}$ where ${J}_{\lambda}$ attains its minimal value, $\displaystyle {m_{\lambda}=J_{\lambda}(u_*)}$.
     \end{itemize}
     
\end{lem}

\begin{proof}

\begin{itemize}
    \item[(i)] We will show that $\mathcal{N}_{\lambda}\ne \emptyset$. Indeed,
since $\lambda>\Lambda_{q}$, it follows from  \eqref{Lambda_q}  that there exists $v_{0} \in W \backslash\{0\}$ such that $\|v_0\|_{c_2,q}^q<\lambda K_q(v_{0})$. 
We claim that there exists $\tau>0$, such that $\tau v_{0} \in \mathcal{N}_{\lambda}$.  Observe that  $\tau v_{0} \in \mathcal{N}_{\lambda}$ means  that 
\begin{equation}\label{tau_v_1}
    \tau^{p} \|v_0\|_{c_1,p}^p+\tau^{q} \|v_0 \|_{c_2,q}^q=\lambda \tau^{q}K_q(v_{0}).
\end{equation}
Dividing  by $\tau^{p},$ equation \eqref{tau_v_1} is equivalent to 
\begin{equation}\label{tau_v_2}
    \|v_0 \|_{c_1,p}^p+\tau^{q-p} \|v_0 \|_{c_2,q}^q=\lambda \tau^{q-p}K_q(v_{0}).\end{equation}

Taking into account that $\lambda K_q(v_{0})-\|v_0\|_{c_2,q}^q>0$, we see that equation \eqref{tau_v_2} can be solved for $\tau$, with 
\begin{equation}\label{tau_v_3}
\tau=\left(\frac{\|v_0\|_{c_1,p}^p}{\lambda K_q(v_{0})-\|v_0\|_{c_2,q}^q}\right)^{\frac{1}{q-p}}
,
\end{equation}
which is positive as $\|v_{0}\|_{c_1,p}>0$. Therefore, choosing  $\tau$ as in \eqref{tau_v_3}, we have $\tau v_{0} \in \mathcal{N}_{\lambda}$.
\item[(ii)]  Let 
$\left(u_{n}\right) \subseteq \mathcal{N}_{\lambda}$ be a minimizing sequence for $J_{\lambda}$. We claim that $(u_{n}) $  is bounded in $W$. Assume to the contrary that $\left(u_{n}\right)$ is unbounded in $W$. Then there exists a subsequence, again denoted $\left(u_{n}\right)$, such that $\left\|u_{n}\right\|_{c_2,q} \rightarrow \infty$. 
Define $v_{n}=u_{n} /\left\|u_{n}\right\|_{c_2,q}.$ Then $\left\|v_{n}\right\|_{c_2,q}=1. $ Since $W$ is reflexive,  one has (by going to a subsequence relabeled $(v_n)$, if need be) that there is $v_{0} \in W$ such that $v_{n} \rightharpoonup v_{0}$ in $W$. Due to the compact embedding of $W^{1,q}(\Omega)$ in $ L^q(\Omega)$ and the continuous embedding of  $L^q(\Omega)$ into $L^p(\Omega)$, (for a subsequence similarly relabeled) $v_{n} \rightharpoonup v_{0}$ in $W^{1, p}(\Omega)$ and  $v_{n} \rightarrow v_{0}$ in $L^{q}(\Omega)$  and $v_{n} \rightarrow v_{0}$ in $L^{p}(\Omega)$.  From \eqref{J_rest_eta} and \eqref{Jlambda}, we have
\begin{equation}\label{Nehari_0}
    \lim_{n\to \infty}\frac{q-p}{q p} \|u_n\|_{c_1,p}^p= m_{\lambda}.
\end{equation}
Since $\displaystyle \lim_{n\to \infty} \|u_n\|_{c_2, q}=\infty$, then from \eqref{Nehari_0} we get 
\begin{equation*}
\lim_{n\to \infty}\frac{q-p}{q p} \|v_n\|_{c_1,p}^p=\lim_{n\to \infty}\frac{\frac{q-p}{q p} \|u_n\|_{c_1,p}^p} {\left\|u_{n}\right\|_{c_2,q}^p} = 0.
\end{equation*}
This implies $ v_n\to 0$ in $W^{1,p}(\Omega)$ as $n\to \infty$. By the continuous embedding of $W^{1,p}(\Omega)$ in $L^p(\Omega)$, it follows that $v_n \to 0$ in $L^p(\Omega)$ as $n\to \infty$. Hence, $v_0(x)=0$ a.e. $x\in \Omega$, which implies $\|v_{n}\|_{q}\to 0$  as $n\to \infty$.
Using the fact that  $u_n \in \mathcal{N}_\lambda$, we have that  
\begin{equation}\label{Nehari_1}0 \leqslant \|u_n \|_{c_1,p}^p=\lambda K_q(u_{n})-\|u_n \|_{c_2,q}^q , \qquad \forall n  \geqslant 1.  \end{equation}
Dividing \eqref{Nehari_1} by $\|u_n\|_{c_2,q}^q$ and taking into account  \eqref{bound-K_q}, we get
\begin{equation}\label{vn}
0 \leqslant \frac{\|u_n \|_{c_1,p}^p}{\|u_n\|_{c_2,q}^q}=\lambda C_1 \|v_{n}\|_{q}^{q}-1 ,\ \ \forall n \geqslant 1 .
\end{equation}
Taking the limit as  $n\to\infty$ in \eqref{vn}  we get $0\leqslant -1.$ This is a contradiction. Thus, $\left(u_{n}\right)$ is bounded in $W$.

\item[(iii)] 
Let us show that $m_{\lambda} >0$. Assume by contradiction that $m_{\lambda}=0$ and let $\left(u_{n}\right) \subseteq \mathcal{N}_{\lambda}$ be a minimizing sequence for $J_{\lambda}$. From $\mbox{(ii)}$  $\left(u_{n}\right)$ is bounded in $W$. 
Since $W$ is reflexive, there exists a subsequence, relabeled again $(u_n)$, such that  $u_{n} \rightharpoonup u_{0}$  in $W$, $u_n\rightharpoonup u_{0}$ in $W^{1, p}(\Omega)$ and $u_{n} \rightarrow u_{0}$ in both $L^{q}(\Omega) $ and  $L^{p}(\Omega)$ because of  the compact embedding of $W^{1,q}(\Omega)$ in $ L^q(\Omega)$ and the continuous embedding of  $L^q(\Omega)$ into $L^p(\Omega)$.
From \eqref{J_rest_eta}, \eqref{Jlambda}  and the fact that $m_\lambda=0$, we have $\|u_n \|_{c_1,p}^p\rightarrow 0.$ Therefore $u_n\to 0 \mbox{ in } L^p({\Omega})$. Hence,  by the uniqueness of limit,  $u_{0} \equiv 0.$
Define $w_{n}=u_{n} /\left\|u_{n}\right\|_{ q}, n \geq 1$.  Observe  $(w_n)$ is bounded in $L^q(\Omega)$ and $\|w_n\|_q=1$ for all $n$. We assert that $\left(w_{n}\right)$ is also bounded in $W$.
Indeed, since $\left(u_{n}\right) \subseteq \mathcal{N}_{\lambda}$, then it follows from \eqref{def_Nehari}  that   $u_n\neq 0$ for each $n$ and $\|u_n\|_{c_1,p}^p+\|u_n \|_{c_2,q}^q=\lambda K_q(u_n)$.  Dividing by $\|u_n\|_q^q$ and taking into account \eqref{bound-K_q}, we obtain  
\begin{equation}\label{w_necu}
\frac{\|w_n \|_{c_1,p}^p}{\left\|u_{n}\right\|_{q}^{q-p}}+\|w_n \|_{c_2,q}^q=\lambda \frac{K_q(u_n)}{\|u_n\|_q^q} \leqslant \lambda C_1.
\end{equation}
Therefore, 
\begin{equation*}\label{w_bound}
    \|w_n \|_{c_2,q}^q\leqslant\lambda C_1,
\end{equation*}
which implies that $\left(w_{n}\right)$ is bounded in $W$. By the reflexivity of $W$, there exists a subsequence, relabeled again $(w_n)$, such that  $w_{n} \rightharpoonup w_{0}$  in $W$, $w_{n} \rightharpoonup w_{0}$ in $W^{1, p}(\Omega)$  and $w_{n} \rightarrow w_{0}$ in $L^{q}(\Omega).$   
From \eqref{w_necu}, we have 
\begin{equation}\label{w_eq}
    \|w_n \|_{c_1,p}^p\leqslant \left\|u_{n}\right\|_{q}^{q-p}(\lambda C_1- \|w_n \|_{c_2,q}^q).
\end{equation}

Using the fact $ (w_n )$ is bounded in $W$,  $u_{n} \rightarrow 0$ in $L^{q} (\Omega)$ and $q-p>0$, we have that \eqref{w_eq} leads to  $\|w_n \|_{c_1,p}\rightarrow 0.$ Then $w_n\rightarrow 0$ in $W^{1,p}(\Omega).$  By the uniqueness of limit, $w_0=0$. Therefore, $w_{n} \rightarrow 0$ in $L^{q}(\Omega)$. Hence, $\|w_n\|_q\to 0$ and $n\to \infty$.  This contradicts the fact that $\|w_n\|_q=1$ for all $n$.
\vspace{0.3cm}

\item[(iv)] Let us show that there exists $u_{*} \in \mathcal{N}_{\lambda}$ such that $J_{\lambda}\left(u_{*}\right)=m_{\lambda}$.\\
Let $\left(u_{n}\right) \subseteq \mathcal{N}_{\lambda}$ be a minimizing sequence, that is, $J_{\lambda}\left(u_{n}\right)=\frac{q-p}{pq}\|u_n \|_{c_1,p}^p \rightarrow m_{\lambda}$ as $n\to \infty$. Then \begin{equation}\label{K_p}
    \lim_{n\to\infty}\|u_n \|_{c_1,p}^p=\frac{qp}{q-p}m_{\lambda}.
\end{equation}
By $\mbox{(ii)}$ $\left(u_{n}\right)$ is bounded in $W$,  and therefore there exists a subsequence $\left(u_{n}\right)$ such that $u_{n} \rightharpoonup u_{*} \mbox{ in } W$ and so, $u_{n} \rightharpoonup u_{*} \in W^{1,p}(\Omega)$ and $u_{n} \to u_{*}$ in $L^{q}(\Omega)$ and $L^{p}( \Omega)$.

It follows from the properties of weak convergence  that
\begin{equation}\label{kp}
    \|u_* \|_{c_1,p}^p \leqslant \liminf_{n\to \infty} \|u_n \|_{c_1,p}^p
\end{equation}
and 
\begin{equation}\label{kq}
    \|u_*\|_{c_2,q}^q\leqslant \liminf_{n\to \infty} \|u_n \|_{c_2,q}^q.
\end{equation}
Since $(u_n)\subseteq \mathcal{N}_{\lambda}$, from the continuity of $K_q$ (see Proposition \ref{prop-cu}) it follows that 
\begin{align*}
    \lambda K_q(u_{*})&= \lambda \lim_{n\to\infty} K_q(u_n)=  \liminf_{n\to \infty}( \|u_n \|_{c_1,p}^p+\|u_n \|_{c_2,q}^q).
    \end{align*}

By \eqref{kp} and \eqref{kq} and the properties of the liminf, we have 

\begin{equation}\label{deskpkq}
    \|u_*\|_{c_1,p}^p+\|u_* \|_{c_2,q}^q\leqslant  \lambda K_q(u_{*}) .
\end{equation}
We assert that $u_{*}\neq 0.$ Assume to the contrary that $u_{*}= 0.$
Using the fact that $u_n\in \mathcal{N}_{\lambda},$ we have that $\|u_n \|_{c_2,q}^q=\lambda K_q(u_n)-\|u_n \|_{c_1,p}^p.$ Since $\lim_{n\to\infty}K_q(u_n)=K_q(0)=0$, then by  \eqref{K_p} we get \[\lim_{n\to\infty} \|u_n \|_{c_2,q}^q=\lambda\lim_{n\to\infty} K_q(u_n)-\lim_{n\to\infty} \|u_n \|_{c_1,p}^p=-\frac{qp}{q-p}m_{\lambda}<0.\]
This leads to a contradiction. Thus, $u_{*}\neq 0.$\\
Now, let us prove that  $\|u_*\|_{c_1,p}^p+\|u_* \|_{c_2,q}^q=\lambda K_q(u_{*})$. Assume by contradiction that the equality does not hold. Then by \eqref{deskpkq}, 
\begin{equation}\label{hpc}
\|u_*\|_{c_1,p}^p+\|u_* \|_{c_2,q}^q<\lambda K_q(u_{*}).
\end{equation}
We claim that there is a constant $\tau \in(0,1)$ such that $\tau u_{*} \in \mathcal{N}_{\lambda}$. To see this,  define a continuous function $g:(0, \infty) \rightarrow \mathbb{R}$ by
$
 g(t):=t^{p-q}\|u_*\|_{c_1,p}^p+\|u_* \|_{c_2,q}^q-\lambda K_q(u_{*}).
$

As $\|u_*\|_{c_1,p}^p>0$ and $p<q$, we have $\lim_{t \rightarrow 0^{+}}g(t)=\infty$.  By \eqref{hpc} we notice that $g(1)<0$. Therefore, there exists $\tau \in(0,1)$ such that $g(\tau)=0$ which implies $\tau u_{*} \in \mathcal{N}_{\lambda}$.
Since $u_{n} \rightharpoonup u_{*} \in W^{1,p}(\Omega)$, we have 
\begin{equation}\label{semi}
    \frac{q-p}{pq}\|u_*\|_{c_1,p}^p \leqslant \liminf_{n\to \infty} \frac{q-p}{pq}\|u_n\|_{c_1,p}^p .
\end{equation}
By \eqref{J_rest_eta}, \eqref{Jlambda}, \eqref{semi} and the fact that $\tau\in (0,1)$ we obtain
$$
m_{\lambda} \leqslant J_{\lambda}\left(\tau u_{*}\right)
\leqslant \tau^{p} \lim _{n \rightarrow \infty} J_{\lambda}\left(u_{n}\right)=\tau^{p} m_{\lambda}<m_{\lambda}.
$$
This is a contradiction. So, 
\begin{equation*}\label{deskp}
\|u_*\|_{c_1,p}^p+\|u_* \|_{c_2,q}^q=\lambda K_q(u_{*}).
\end{equation*}
This implies that  $u_{*}\in \mathcal{N}_{\lambda}$.  By \eqref{Jlambda} and \eqref{semi}, we have

$$m_{\lambda} \leqslant J_{\lambda}\left(u_{*}\right)= \frac{q-p}{q p} \|u_*\|_{c_1,p}^p \leqslant \liminf_{n\to \infty} \frac{q-p}{pq}\|u_n\|_{c_1,p}^p = m_{\lambda}.
$$
Thus $J_{\lambda}\left(u_{*}\right)=m_{\lambda}$. This completes the proof.
\end{itemize}
\end{proof}

\end{document}
